\chapter{线程}

\section{线程定义}
\concept{为什么引入线程}
\begin{points}
  \item 引入进程是为了描述和实现多道程序并发执行。
  \item 引入线程是为了减少并发执行的时空开销，让同一应用内部的多个任务能更轻量地并发。
  \item 早期进程同时承担“资源拥有单位”和“调度分派单位”两个角色，创建和切换开销较大。
  \item 线程把二者分开：进程更多作为资源拥有单位，线程作为 CPU 调度的基本执行单位。
\end{points}
\plain{线程是进程里面更轻的执行流，同进程线程共享资源，但各自有栈、寄存器和 PC。}
\exam{老师明确点过线程代码题和“为什么引入线程”，常考把进程的两个属性拆开来答。}


\concept{线程的定义}
\begin{points}
  \item 线程是进程内的一个执行单元，是进程内的可调度实体。
  \item 一个线程由线程 ID、程序计数器、寄存器集合和栈组成。
  \item 同一进程内的线程共享代码段、数据段、打开文件、地址空间等资源。
  \item 线程是“轻量级进程”，但不是独立资源拥有者。
\end{points}
\plain{定义题抓两半写：线程自己只带执行现场，资源主体仍挂在所属进程名下。}
\exam{常考线程由哪些部分组成，以及哪些资源是线程私有、哪些是进程内共享。}


\concept{线程私有资源与共享资源}
\begin{points}
  \item 线程私有的典型内容包括线程 ID、程序计数器、寄存器集合和执行栈。
  \item 同一进程内线程共享代码段、数据段、堆、打开文件及多数操作系统资源。
  \item 正因为共享地址空间，线程间通信方便，但一个线程写坏共享数据也更容易影响同进程其他线程。
\end{points}
\plain{记忆法：凡是“执行到哪一步”的现场通常归线程，凡是“这整个程序拥有什么”的资源通常归进程。}
\exam{常考判断题：栈、寄存器、PC 是否共享；代码段、全局数据、打开文件是否共享。}


\concept{进程与线程区别}
\begin{points}
  \item 进程是资源分配的基本单位，线程是 CPU 调度的基本单位。
  \item 同一进程内线程共享地址空间，线程之间通信方便但相互影响更大。
  \item 不同进程地址空间隔离，通信开销较大但保护性更好。
  \item 线程创建、撤销、切换通常比进程轻量，因为不需要切换完整地址空间和大量资源。
\end{points}
\plain{这题别答成“线程更小的进程”就结束，必须把“资源归进程、调度归线程、共享地址空间”三点写出来。}
\exam{常考进程和线程在调度单位、资源拥有、通信开销、切换开销上的对比。}


\concept{多线程优点}
\begin{points}
  \item 响应性更好：例如界面线程响应用户，后台线程执行耗时任务。
  \item 资源共享方便：同一进程内线程天然共享内存和文件。
  \item 经济性更好：创建和切换开销较小。
  \item 可伸缩性更好：多核系统中多个线程可并行执行。
\end{points}
\plain{PPT 这里强调四个关键词：响应度、共享、经济性、可伸缩性，答简答题按这四点展开最稳。}
\exam{常考“多线程有哪些优势”，最好按响应性/共享/经济性/可伸缩性分点写。}


\concept{并发与并行}
\begin{points}
  \item 并发 concurrency 指多个任务在同一时间段内交替推进。
  \item 并行 parallelism 指多个任务在同一时刻真正同时执行，通常依赖多核处理器。
  \item 单核系统通常体现为并发；多核系统既能并发，也能并行。
\end{points}
\plain{并发强调“都在推进”，并行强调“真的同时跑”。}
\exam{常考辨析题：单核只能并发不能真正并行，多核才可能并行。}


\concept{Amdahl 定律}
\begin{points}
  \item 设程序中串行部分比例为 $S$，并行部分比例为 $1-S$，处理器核心数为 $N$。
  \item 理想情况下，总执行时间满足：$T_{new}=S\times T_{old}+(1-S)\times \dfrac{T_{old}}{N}$。
  \item 当 $N$ 趋于无穷大时，加速比上限接近 $1/S$。
  \item 因此串行部分越大，多核带来的性能提升上限越低。
\end{points}
\plain{再多的核也救不了那段必须串行跑的代码，串行部分像木桶最短板。}
\exam{常考给定串行比例和核数算理论加速比，或说明为什么加核数后收益会递减。}


\concept{多核编程的挑战}
\begin{points}
  \item 任务分解：识别哪些子任务能够独立并发执行。
  \item 负载均衡：尽量让多个核心工作量接近，避免一核忙一核闲。
  \item 数据划分：让不同核心处理不同数据块，减少争抢。
  \item 数据依赖与同步：正确处理共享数据和先后约束。
  \item 测试调试：并发程序的错误具有时序相关性，更难复现。
\end{points}
\plain{多核难点不在“开几个线程”，而在“怎么拆、怎么同步、怎么别互相拖后腿”。}
\exam{常考简答：为什么多核程序比单线程程序更难设计和调试。}


\concept{数据并行与任务并行}
\begin{points}
  \item \term{数据并行}：把同一种操作分配到不同数据块上，由多个核心并行处理。
  \item \term{任务并行}：把不同功能的任务分给不同执行流，各自执行不同操作。
  \item 数据并行更强调“同算法处理不同数据”；任务并行更强调“不同任务分工协作”。
\end{points}
\plain{看题目先分辨是在“把一堆数据切开一起算”，还是“把不同工作拆给不同线程做”，前者是数据并行，后者是任务并行。}
\exam{PPT 在多核编程挑战后专门给了这两个概念，常考两者区别或让你判断一个场景属于哪一类并行。}


\section{线程实现}
\concept{用户级线程}
\begin{points}
  \item 用户级线程由用户态线程库管理，内核不知道线程存在。
  \item 优点：创建和切换无需进入内核，速度快；调度策略可由应用定制。
  \item 缺点：一个线程执行阻塞系统调用，可能导致整个进程阻塞；无法真正利用多核并行。
\end{points}
\plain{用户级线程的优势几乎都来自“不进内核”，它的主要短板也正因为“内核根本看不见它”。}
\exam{常考用户级线程为什么快、为什么一阻塞可能整个进程都受影响。}


\concept{内核级线程}
\begin{points}
  \item 内核级线程由操作系统内核管理，调度器直接调度线程。
  \item 优点：一个线程阻塞不影响同进程其他线程；多线程可在多核上并行。
  \item 缺点：线程创建、切换、同步需要内核参与，开销比用户级线程大。
\end{points}
\plain{内核级线程的关键点是“内核真能看见并分别调度每个线程”，所以一个线程卡住时别的线程还能继续。}
\exam{常考为什么内核级线程的阻塞不会拖死同进程其他线程，以及它为什么更能利用多核。}


\concept{为什么内核级线程阻塞不影响其他线程}
\begin{points}
  \item 在内核级线程实现中，处理机分配对象是线程而不是整个进程。
  \item 某个线程因系统调用或 I/O 阻塞时，内核仍能把 CPU 分给同进程的其他就绪线程。
  \item 因此“一个线程阻塞，整个进程全停住”的现象通常发生在用户级线程或多对一模型中，而不是内核级线程中。
\end{points}
\plain{因为内核调度表里登记的是“线程”，不是“这个进程整体”，所以它能跳过被阻塞的那个线程。}
\exam{常考 Why 题：为什么一个内核级线程阻塞后，同进程其他线程仍可运行。}


\concept{用户级线程为什么阻塞整个进程}
\begin{points}
  \item 操作系统看不见用户级线程，内核调度对象仍是整个进程。
  \item 若某个用户级线程执行阻塞系统调用，内核只知道“这个进程阻塞了”。
  \item 因而同一进程中的其他用户级线程也失去运行机会，直到该系统调用返回。
\end{points}
\plain{内核眼里根本没有这些用户线程，所以一个线程卡住时，内核会把整个进程一起挂起。}
\exam{常考用户级线程的典型缺点，尤其与多对一模型一起考。}


\concept{纯用户级线程的调度与抢占限制}
\begin{points}
  \item 纯用户级线程通常由线程库在用户空间完成调度，内核调度对象依然是进程。
  \item 若没有专门机制配合，线程之间通常不具备像内核调度那样的强制抢占能力。
  \item 因此纯用户级线程常依赖线程主动让出处理器，例如执行结束、调用调度函数，或显式调用 \code{thread\_yield()}。
\end{points}
\plain{纯用户级线程更像“进程自己在内部排班”，不主动交棒时，别的用户线程往往很难硬抢过来。}
\exam{常考讨论题：纯用户级线程是否支持抢占、如何切换、为何常需要主动让出 CPU。}


\concept{多线程模型}
\begin{points}
  \item 多对一模型：多个用户线程映射到一个内核线程，管理轻量但阻塞问题明显。
  \item 一对一模型：每个用户线程映射到一个内核线程，并行能力强但内核线程数量受限。
  \item 多对多模型：多个用户线程映射到多个内核线程，兼顾灵活性与并行能力，但实现复杂。
\end{points}
\plain{模型题核心不是背名字，而是看“用户线程和内核线程怎么映射”，再推导阻塞、并行和开销结果。}
\exam{常考多对一/一对一/多对多各自的阻塞特性、能否利用多核、实现代价。}


\concept{多对一模型}
\begin{points}
  \item 多个用户级线程映射到一个内核线程，线程管理主要由用户空间线程库完成。
  \item 优点是创建和管理轻量。
  \item 缺点是一个线程一旦执行阻塞系统调用，整个进程都会阻塞；同时任一时刻只有一个线程能进入内核，无法真正利用多核。
  \item 典型例子包括早期 Solaris Green Threads、GNU Portable Threads。
\end{points}
\plain{多对一最像“很多演员共用一个内核身份”，谁一旦在内核这层卡住，整组人都得等。}
\exam{常考多对一模型为什么阻塞问题最严重，以及为什么它无法发挥多核优势。}


\concept{一对一模型}
\begin{points}
  \item 每个用户线程映射到一个内核线程。
  \item 一个线程阻塞时，其他线程仍可被内核继续调度。
  \item 它支持多线程在多核系统上并发甚至并行运行。
  \item 代价是大量创建内核线程会带来更高系统开销，因此系统通常会限制线程数量。
\end{points}
\plain{一对一的好处最直接，但代价也最直接: 每多一个用户线程，内核里就真多一个调度实体。}
\exam{常考一对一模型相对多对一为什么更适合现代多核系统，以及它的主要成本。}


\concept{多对多模型}
\begin{points}
  \item 多个用户线程映射到少量或相等数量的内核线程。
  \item 用户可创建较多线程，而内核线程可在多处理器系统上并发运行。
  \item 当某个线程执行阻塞系统调用时，内核仍可调度其他内核线程继续推进。
  \item 其变种二层模型还允许部分线程和内核线程做一对一绑定。
\end{points}
\plain{多对多是在用户灵活性和内核并行能力之间折中: 用户线程很多，但不必每个都绑一个内核线程。}
\exam{常考多对多模型如何兼顾用户级线程灵活性和内核级线程并行能力。}


\concept{线程库}
\begin{points}
  \item 线程库向程序员提供创建、管理、同步线程的 API。
  \item 常见线程库包括 POSIX Pthreads、Windows threads、Java threads。
  \item 线程库可以完全在用户态实现，也可以由操作系统提供内核级线程支持。
\end{points}
\plain{线程库是程序员看到的 API 外壳，真正要区分的是“库函数只在用户态跑完”，还是“调用后会进内核”。}
\exam{常考线程库的作用，以及它和用户级线程/内核级线程实现之间的关系。}


\concept{线程库的两种实现方式}
\begin{points}
  \item \term{用户空间线程库}：代码和数据结构都在用户空间，调用 API 后通常只是过程调用，不触发系统调用。
  \item \term{内核支持的线程库}：线程控制依赖操作系统内核，调用 API 后往往会触发系统调用，由内核完成创建、调度和管理。
  \item 前者快但内核不可见，后者功能强且支持真正并行，但模式切换开销更大。
\end{points}
\plain{同样叫“线程库”，实现差别可能很大，考试常让你顺着“是否进内核”推出阻塞、并行能力和开销。}
\exam{常考线程库的两种实现方式，以及它们和用户态/内核态支持的关系。}


\concept{常见线程库与 API}
\begin{points}
  \item Pthreads 是 POSIX 线程标准，常见接口包括 \code{pthread\_create}、\code{pthread\_join}、\code{pthread\_exit}、\code{pthread\_cancel}。
  \item Win32 线程库常见接口包括 \code{CreateThread}、\code{ExitThread}、\code{WaitForSingleObject}。
  \item 这类 API 是代码题高频背景，要能看懂“创建几个线程、谁等待谁、哪些资源共享”。
\end{points}
\plain{考试不要求你背完整手册，但看到这些函数名要立刻知道它们在干什么。}
\exam{老师录音明确点过线程代码题，这里常和输出分析、共享变量分析一起出大题。}


\concept{\code{pthread\_create} 与 \code{pthread\_join}}
\begin{points}
  \item \code{pthread\_create} 用于创建一个新线程，指定线程标识、属性、线程函数和传入参数。
  \item 新线程与创建者属于同一进程，因此通常共享代码段、全局数据、堆和打开文件等资源。
  \item \code{pthread\_join} 用于等待指定线程结束，常用于保证主线程在子线程执行完后再继续。
  \item 代码题里若子线程修改全局变量，主线程在 \code{pthread\_join} 之后通常能看到修改结果；这与 \code{fork()} 后父子进程普通变量互不影响形成对比。
\end{points}
\plain{\code{pthread\_create} 重点看“新线程共享哪些资源”，\code{pthread\_join} 重点看“主线程是不是等它干完再往下执行”。}
\exam{线程代码题高频考点就是创建了几个线程、是否等待、共享变量最终值是多少，以及和 \code{fork()} 代码的差别。}


\section{操作系统中的线程例子}
\concept{Windows 线程}
\begin{points}
  \item Windows 采用一对一模型，每个用户线程对应一个内核线程。
  \item 一个线程通常包含线程 ID、寄存器组、用户栈、内核栈和私有存储区。
  \item 线程相关数据结构常见有 ETHREAD、KTHREAD、TEB。
\end{points}
\plain{Windows 线程是“每建一个用户线程，内核里真有一个对应实体”。}
\exam{常考操作系统例子时，重点不是背缩写全称，而是知道 Windows 基本上是一对一模型。}


\concept{Linux 线程}
\begin{points}
  \item Linux 中常把线程和进程都抽象成 task，线程创建常通过 \code{clone()} 实现。
  \item \code{fork()} 创建的新执行流通常不共享地址空间和大多数资源。
  \item \code{clone()} 可按标志位决定是否共享地址空间、文件表等资源，因此可用来实现线程。
  \item 在类 Unix 系统中，Pthreads 通常基于内核级线程的一对一模型实现。
  \item 区分 \code{fork} 和 \code{clone} 的关键在于“是否共享资源与地址空间”。
\end{points}
\plain{Linux 不是单独造一个完全不同的“线程生物种类”，而是靠共享资源程度不同来区分。}
\exam{常考代码题或概念题：为什么 \code{fork} 后普通变量互不影响，而线程共享全局数据。}


\concept{可重入代码}
\begin{points}
  \item 可重入代码指在被多个执行流同时调用或被中断后再次进入时，仍能得到正确结果的代码。
  \item 它通常不依赖不可保护的全局变量，不在函数内部保留共享静态状态，或对共享状态进行同步保护。
  \item 多线程、信号处理和中断处理程序常要求代码具备可重入性。
  \item 直觉：可重入代码像“每个人带自己的草稿纸”，不会抢同一张草稿纸写中间结果。
\end{points}
\plain{可重入代码像“公共只读函数”，多个线程同时进来也不会互相污染。}
\exam{常考可重入代码特征：不修改自身代码、不依赖全局可变状态或对其加保护。}
